% =============================================================================
\section{Core Framework: Influence Profiles and Effective Context Length}\label{sec:framework}
% =============================================================================

\subsection{Single-position influence energy}

Fix reference position $r$. For each position $i$, define a single-site perturbation $X^{(i)}\sim \Pi_{\{i\}}(\cdot\mid X)$ and define the \emph{influence energy}
\begin{equation}\label{eq:single_influence}
\cI(i;r) \;:=\; \E\!\left[d_\cZ\!\left(f(X),\, f(X^{(i)})\right)\right].
\end{equation}
This measures the expected embedding change induced by perturbing position $i$ while holding other positions fixed.

\begin{remark}[Dependence on perturbation choice]
$\cI(i;r)$ depends on $\Pi_{\{i\}}$. Random substitution, dinucleotide shuffle, and conditional resampling yield different magnitudes. Our goal is not to remove this dependence but to formalize it and recommend reporting ECL as a function of perturbation type.
\end{remark}

\subsection{Normalized influence profile}

Define the \emph{normalized influence profile} as a probability distribution over positions:
\begin{equation}\label{eq:normalized_influence}
\bar{\cI}(i;r) \;:=\; \frac{\cI(i;r)}{\cI_{\mathrm{tot}}(r)}, \qquad
\cI_{\mathrm{tot}}(r) := \sum_{i=1}^L \cI(i;r),
\end{equation}
assuming $\cI_{\mathrm{tot}}(r) > 0$. The normalized profile $\bar{\cI}(\cdot;r)$ is a probability mass function on $\{1,\dots,L\}$, interpretable as the probability that a random perturbation effect originates at position $i$.

\subsection{Distance-binned influence profile}

Define the binned influence profile as a function of distance $d\ge 0$:
\begin{equation}\label{eq:distance_profile}
\cI(d;r) := \frac{1}{|\{i:|i-r|=d\}|}\sum_{i:|i-r|=d} \cI(i;r),
\end{equation}
with the convention that empty sums are zero near boundaries.

\subsection{Cumulative influence and effective context length}

Define cumulative influence within radius $\ell$:
\begin{equation}\label{eq:cum_influence}
\cI_{\le \ell}(r) := \sum_{i:|i-r|\le \ell} \cI(i;r)
= \sum_{d=0}^{\ell} |\{i:|i-r|=d\}|\, \cI(d;r).
\end{equation}

\begin{definition}[Perturbation--variance effective context length]\label{def:ecl}
Fix $\beta\in(0,1]$. The \emph{effective context length at level $\beta$} is
\begin{equation}\label{eq:ecl_def}
\ECL_\beta(r) \;:=\; \min\!\left\{\ell\in\{0,\dots,L\}:\; \cI_{\le \ell}(r) \ge \beta\, \cI_{\mathrm{tot}}(r)\right\}.
\end{equation}
\end{definition}

This is a population-level definition. For finite datasets, we estimate $\cI(i;r)$ empirically and plug in.

\subsection{Effective Context Profile (ECP)}

\begin{definition}[Effective context profile]\label{def:ecp}
The \emph{effective context profile} is the function $\ECP_r:(0,1] \to \{0,\dots,L\}$ defined by $\ECP_r(\beta) := \ECL_\beta(r)$. It is the complete non-decreasing curve characterizing context utilization, strictly more informative than any single $\ECL_\beta$.
\end{definition}

We define a scalar summary:
\begin{definition}[Area under the ECP]\label{def:aecp}
The \emph{area under the effective context profile} is
\begin{equation}\label{eq:aecp}
\AECP(r) := \int_0^1 \ECL_\beta(r)\, d\beta.
\end{equation}
\end{definition}

\begin{proposition}[AECP equals mean influence distance]\label{prop:aecp}
$\AECP(r) = \sum_{i=1}^L |i-r| \cdot \bar{\cI}(i;r)$.
\end{proposition}
\begin{proof}
Let $D$ denote the random distance from $r$ with $\Pr(D = d) = \sum_{i:|i-r|=d}\bar{\cI}(i;r)$, so that $D$ has cumulative distribution function $F(d) = \cI_{\le d}(r)/\cI_{\mathrm{tot}}(r)$ and quantile function $Q(\beta) = \min\{d : F(d) \ge \beta\}$. By definition, $\ECL_\beta(r) = Q(\beta)$. The discrete quantile identity
\[
\int_0^1 Q(\beta)\,d\beta \;=\; \sum_{d \ge 0} d \cdot \Pr(D = d) \;=\; \E[D]
\]
holds because, for the step quantile function of a discrete distribution, $\int_0^1 Q(\beta)\,d\beta$ telescopes to $\sum_d d \cdot (F(d) - F(d-1)) = \sum_d d \cdot \Pr(D=d)$. Substituting the explicit form of $\Pr(D = d)$ yields $\E[D] = \sum_{i=1}^L |i-r| \bar{\cI}(i;r)$.
\end{proof}

\subsection{Directional ECL}

Genomic regulation is directionally asymmetric: promoters lie upstream, $3'$-UTR effects downstream. We adopt the convention that ``upstream'' refers to indices $i < r$ (smaller coordinate, by convention closer to the $5'$ end on the plus strand) and ``downstream'' to indices $i > r$, with the understanding that strand orientation must be tracked when comparing across loci. Define upstream and downstream influence:
\begin{align}
\cI^-(d;r) &:= \cI(r-d;r), \qquad d>0, \label{eq:upstream}\\
\cI^+(d;r) &:= \cI(r+d;r), \qquad d>0, \label{eq:downstream}
\end{align}
and corresponding directional ECLs:
\begin{equation}\label{eq:directional_ecl}
\ECL_\beta^-(r) := \min\!\left\{\ell : \sum_{d=1}^{\ell} \cI^-(d;r) \ge \beta \sum_{d=1}^{L} \cI^-(d;r)\right\},
\end{equation}
and analogously for $\ECL_\beta^+(r)$. Asymmetry is quantified by the ratio $\ECL_\beta^+(r)/\ECL_\beta^-(r)$.

\subsection{Effective Context Dimension (ECD)}

\begin{definition}[Effective context dimension]\label{def:ecd}
The \emph{effective context dimension} is the exponential of the entropy of the normalized influence profile:
\begin{equation}\label{eq:ecd}
\ECD(r) := \exp\!\left(-\sum_{i=1}^L \bar{\cI}(i;r) \log \bar{\cI}(i;r)\right).
\end{equation}
\end{definition}

ECD measures the effective number of positions contributing to the embedding, analogous to the participation ratio in physics. A model concentrating influence on $k$ positions has $\ECD \approx k$; one spreading influence uniformly over $m$ positions has $\ECD = m$. ECD is invariant to the total scale of $\cI$ and is bounded: $1 \leq \ECD(r) \leq \abs{\mathrm{supp}(\bar{\cI}(\cdot;r))} \leq L$, with the lower bound attained when influence concentrates on a single position. The fully degenerate case $\cI_{\mathrm{tot}}(r) = 0$ (e.g., when $f$ is constant on the relevant context window) leaves $\bar{\cI}$ undefined; in that case ECD is conventionally set to zero and a flag is reported alongside the estimate.

\subsection{Interaction ECL}

Single-site perturbations miss cooperative effects between distant positions. Define the \emph{pairwise interaction influence}:
\begin{equation}\label{eq:interaction}
\cI_{\mathrm{int}}(i,j;r) := \cI(\{i,j\};r) - \cI(i;r) - \cI(j;r),
\end{equation}
where $\cI(\{i,j\};r) := \E[d_\cZ(f(X), f(X^{(\{i,j\})}))]$ is the joint perturbation influence. Positive $\cI_{\mathrm{int}}$ indicates synergistic (super-additive) effects; negative indicates redundancy. We define the corresponding ECL.

\begin{definition}[Interaction ECL]\label{def:ecl_int}
Fix $\beta\in(0,1]$. The \emph{interaction ECL} at level $\beta$ is the minimal radius capturing $\beta$-fraction of total absolute interaction energy:
\begin{equation}\label{eq:ecl_int_def}
\ECLint_\beta(r) \;:=\; \min\!\left\{\ell\in\{0,\dots,L\} \;:\;
\sum_{\substack{i,j\,:\,\min(|i-r|,|j-r|)\le \ell}} \abs{\cI_{\mathrm{int}}(i,j;r)}
\;\ge\; \beta \sum_{i,j} \abs{\cI_{\mathrm{int}}(i,j;r)}\right\}.
\end{equation}
\end{definition}

\noindent
The absolute value is used in the numerator and denominator so that synergistic and redundant interactions both contribute to the radius, revealing cooperativity between proximal and distal elements regardless of sign.

\subsection{Context Decay Spectroscopy (CDS)}

\begin{definition}[Context decay spectroscopy]\label{def:cds}
Fit the distance-binned influence profile to a mixture of exponentials:
\begin{equation}\label{eq:cds}
\cI(d;r) \;\approx\; \sum_{k=1}^K a_k \, e^{-\lambda_k d},
\end{equation}
where $a_k > 0$ are amplitudes and $\lambda_k > 0$ are decay rates, estimated via non-negative least squares or EM. Each component represents a distinct \emph{context channel}: fast-decaying components ($\lambda_k$ large) capture local motif detection; slow-decaying components ($\lambda_k$ small) capture long-range regulatory integration. The \emph{spectral ECL} at level $\beta$ can be computed analytically from the fitted mixture.
\end{definition}

CDS is more informative than a single ECL number: it reveals whether context utilization is unimodal (one decay scale) or multi-scale (reflecting architecturally distinct processing channels).

\subsection{Window ablation and complementary definitions}

An alternative viewpoint perturbs \emph{all positions outside} a window and measures resulting embedding deviation. Define $S_\ell := \{i:|i-r|>\ell\}$ and:
\begin{equation}\label{eq:out_ablation}
\Delta_{\mathrm{out}}(\ell;r) := \E\!\left[d_\cZ\!\left(f(X),\, f(X^{(\mathrm{out},\ell)})\right)\right],
\end{equation}
where $X^{(\mathrm{out},\ell)}\sim \Pi_{S_\ell}(\cdot\mid X)$. Small $\Delta_{\mathrm{out}}(\ell;r)$ indicates that radius $\ell$ suffices. Under approximate additivity (\cref{prop:tail_add}), $\Delta_{\mathrm{out}}(\ell;r) \approx \sum_{|i-r|>\ell} \cI(i;r)$.

\subsection{Global and task-aware ECL}

Aggregate ECL across loci by $\ECL_\beta := \E_R[\ECL_\beta(R)]$ for a distribution of reference loci $R$ (e.g., transcription start sites). Condition on locus type, chromatin state, or gene class for class-specific ECL.

For a downstream readout $g:\cZ\to\cY_{\mathrm{task}}$, define task-aware influence energies $\cI^{\mathrm{task}}(i;r) := \E[d_{\mathrm{task}}(g(f(X)), g(f(X^{(i)})))]$ with $d_{\mathrm{task}}$ a task loss, yielding task-aware ECL $\ECL^{\mathrm{task}}_\beta(r)$.

\subsection{Genomic Needle-in-a-Haystack (gNIAH)}\label{sec:gniah}

We propose a genomic analogue of the NLP Needle-in-a-Haystack test. Insert a known regulatory motif (e.g., a CTCF binding site or a GATA motif) at varying distances $d$ from the prediction center in a neutral (shuffled) background sequence. Define the \emph{gNIAH sensitivity}:
\begin{equation}\label{eq:gniah}
\gNIAH(d, m) := \E\!\left[d_\cZ\!\big(f(X_{\mathrm{neutral}}),\, f(X_{\mathrm{neutral}}^{+m@d})\big)\right],
\end{equation}
where $X_{\mathrm{neutral}}^{+m@d}$ denotes the neutral sequence with motif $m$ inserted at distance $d$ from the prediction center $r$, and the discrepancy is evaluated at $r$. Equivalently, $\gNIAH(d,m)$ is a task-aware influence energy under an ``insertion'' perturbation kernel that overwrites a contiguous block at distance $d$ from $r$ with the motif $m$, in contrast to the ``replacement'' kernels of \cref{def:perturbation}. The gNIAH profile $d \mapsto \gNIAH(d,m)$ directly measures how far the model can ``see'' a specific regulatory signal, without requiring perturbation of natural sequences.
